\begin{figure}[htbp]
\centering
\resizebox{\textwidth}{!}{%
\begin{tikzpicture}[
    >=Stealth,
    node distance=0.45cm and 0.55cm,
    every node/.style={font=\footnotesize},
    component/.style={
        rectangle,
        rounded corners=4pt,
        draw=black,
        thick,
        minimum width=2.4cm,
        minimum height=1.3cm,
        align=center,
        inner sep=4pt
    },
    seed/.style={component, fill=green!15},
    fuzzer/.style={component, fill=blue!12},
    format/.style={component, fill=red!10},
    execution/.style={component, fill=orange!18},
    logging/.style={component, fill=red!18},
    eval/.style={component, fill=yellow!22},
    output/.style={component, fill=teal!15},
    phase/.style={
        rectangle,
        draw=black!50,
        dashed,
        thick,
        rounded corners=6pt,
        inner sep=10pt
    },
    phaselabel/.style={
        font=\small\bfseries\itshape,
        text=black!70
    },
    arrow/.style={->, thick, draw=black!75}
]

% --- Componentes da Fase de Campanha ---
\node[seed] (seeds) {Seeds /\\Templates};
\node[component, fill=green!22, right=of seeds] (input) {Input\\Generator};
\node[fuzzer, right=of input] (fuzzer) {Fuzzer\\Module\\{\scriptsize\itshape\color{black!60}(× 6 engines)}};
\node[format, right=of fuzzer] (format) {Prompt-Formatting\\Unit};
\node[execution, right=of format] (exec) {Execution\\Module\\{\scriptsize\itshape\color{black!60}(× 3 provedores)}};
\node[logging, right=of exec] (log) {Logging\\Component};

% --- Componentes da Fase de Avaliação ---
\node[eval, below=1.8cm of format] (evalc) {Evaluation\\Component};
\node[output, right=of evalc] (results) {Resultados\\(CSVs)};

% --- Setas da Fase de Campanha ---
\draw[arrow] (seeds) -- (input);
\draw[arrow] (input) -- (fuzzer);
\draw[arrow] (fuzzer) -- (format);
\draw[arrow] (format) -- (exec);
\draw[arrow] (exec) -- (log);

% --- Transição entre fases ---
\draw[arrow] (log.south) -- ++(0,-0.5) -| (evalc.north);

% --- Setas da Fase de Avaliação ---
\draw[arrow] (evalc) -- (results);

% --- Caixa pontilhada Fase de Campanha ---
\begin{scope}[on background layer]
    \node[phase, fit=(seeds)(input)(fuzzer)(format)(exec)(log), inner sep=12pt] (campaign_box) {};
\end{scope}
\node[phaselabel, above=2pt of campaign_box.north west, anchor=south west] {Fase de Campanha (coleta de dados)};

% --- Caixa pontilhada Fase de Avaliação ---
\begin{scope}[on background layer]
    \node[phase, fit=(evalc)(results), inner sep=12pt] (eval_box) {};
\end{scope}
\node[phaselabel, below=2pt of eval_box.south west, anchor=north west] {Fase de Avaliação (aplicação dos oráculos)};

\end{tikzpicture}%
}
\caption{Arquitetura modular do \textit{framework} de \textit{fuzzing} ético, organizada em duas fases independentes: \textbf{Campanha} (coleta de dados via APIs dos LLMs) e \textbf{Avaliação} (aplicação determinística dos oráculos sobre os dados coletados). A separação física dos artefatos entre as fases preserva a imutabilidade dos dados primários e permite a re-execução dos oráculos sem repetição das chamadas de API.}
\label{fig:arquitetura}
\end{figure}